\chapter{Quick Start}
\label{ch:quickstart}

This chapter walks through creating, configuring, and building an OmniLaTeX
document from scratch.  By the end, you will have a complete article with a
figure, a table, and a bibliography entry.

\section{Minimal Document}

The smallest possible OmniLaTeX document requires only five lines:

\begin{minted}{latex}
\documentclass{omnilatex}
\begin{document}
\section{Hello World}
This is a minimal OmniLaTeX document.
\end{document}
\end{minted}

Each line serves a purpose:

\begin{enumerate}
    \item \texttt{\textbackslash documentclass\{omnilatex\}} -- Loads the
          OmniLaTeX class.  With no options specified, all defaults are applied:
          \textsf{doctype=article}, \textsf{language=english}, and the default
          KOMA-Script base class (\texttt{scrartcl}).
    \item \texttt{\textbackslash begin\{document\}} -- Marks the start of the
          document body.  The preamble (everything before this line) is where
          you place configuration, package loading, and metadata.
    \item \texttt{\textbackslash section\{Hello World\}} -- A top-level
          section heading.  OmniLaTeX inherits KOMA-Script's sectioning
          commands, which support optional short titles for the table of
          contents.
    \item The paragraph text -- Rendered in the default Libertinus Serif font
          with \ctanpackage{microtype} optimisations applied automatically.
    \item \texttt{\textbackslash end\{document\}} -- Ends the document body
          and triggers finalisation (writing the table of contents, glossary,
          bibliography, etc.).
\end{enumerate}

Compile with:

\begin{minted}{bash}
latexmk -lualatex main.tex
\end{minted}

\section{Adding Options}

OmniLaTeX accepts key-value options via the \texttt{\textbackslash documentclass}
command.  Options are comma-separated in square brackets.

\subsection{Language}

\begin{minted}{latex}
\documentclass[language=german]{omnilatex}
\end{minted}

This loads \ctanpackage{polyglossia} and sets German as the document language,
activating hyphenation patterns and localised strings (e.g.\ ``Inhaltsverzeichnis''
instead of ``Table of Contents'').  Supported languages include English, German,
French, Japanese, Chinese, Korean, Arabic, and Hebrew.

\subsection{Document Type}

\begin{minted}{latex}
\documentclass[doctype=report]{omnilatex}
\end{minted}

The \textsf{doctype} option selects the document profile.  Each doctype
configures the KOMA-Script base class, page layout, and feature set.
Available doctypes include \textsf{article}, \textsf{report},
\textsf{thesis}, \textsf{book}, \textsf{presentation},
\textsf{poster}, and \textsf{cv}.  The full doctype system is documented in
Chapter~\ref{ch:doctype-system}.

\subsection{Combining Options}

Options can be combined freely:

\begin{minted}{latex}
\documentclass[
  doctype=thesis,
  language=english,
  loadGlossaries,
]{omnilatex}
\end{minted}

The \textsf{loadGlossaries} option activates the
\ctanpackage{glossaries-extra} integration, loading the symbol and
abbreviation management subsystem.

\section{Your First Real Document}

The following example demonstrates a complete OmniLaTeX article with a title,
author, abstract, sections, a figure, a table, and a bibliography entry.

\begin{minted}{latex}
\documentclass[
  doctype=article,
  language=english,
]{omnilatex}

\title{An Analysis of Thermal Conductivity in Nanofluids}
\author{Jane Doe}
\date{\today}

\begin{document}
\maketitle

\begin{abstract}
  This article presents experimental measurements of thermal
  conductivity in water-based alumina nanofluids at volume fractions
  ranging from 0.5\% to 4.0\%.  Results show a non-linear enhancement
  that deviates from the Maxwell model at higher concentrations.
\end{abstract}

\section{Introduction}
Nanofluids -- suspensions of nanometre-sized particles in base fluids
-- have attracted significant interest for heat transfer enhancement
\cite{choi1995nanofluids}.

\section{Methodology}
Experiments were conducted using a transient hot-wire apparatus.  The
test section and measurement procedure are shown in
\Cref{fig:apparatus}.

\begin{figure}[htbp]
  \centering
  \fbox{\parbox{0.6\textwidth}{\centering\vspace{2cm}%
    Schematic of the transient hot-wire apparatus\vspace{2cm}}}
  \caption{Experimental apparatus for thermal conductivity measurement.}
  \label{fig:apparatus}
\end{figure}

\section{Results}
\Cref{tab:results} summarises the measured thermal conductivity ratios.

\begin{table}[htbp]
  \centering
  \caption{Thermal conductivity enhancement at $T = 25\,\si{\degreeCelsius}$}
  \label{tab:results}
  \begin{tabular}{@{} S[table-format=1.1] S[table-format=1.3] @{}}
      \toprule
      {$\phi$ [\%]} & {$k_{\mathrm{nf}} / k_{\mathrm{bf}}$} \\
      \midrule
      0.5 & 1.023 \\
      1.0 & 1.051 \\
      2.0 & 1.098 \\
      4.0 & 1.142 \\
      \bottomrule
  \end{tabular}
\end{table}

\begin{thebibliography}{1}
  \bibitem{choi1995nanofluids}
    S.~U.~S.~Choi, ``Enhancing thermal conductivity of fluids with
    nanoparticles,'' \emph{ASME Publications FED}, vol.~231, pp.~99--105,
    1995.
\end{thebibliography}
\end{document}
\end{minted}

Key observations about this example:

\begin{itemize}
    \item The \texttt{S} column type from \ctanpackage{siunitx} is used for
          aligned numeric columns with automatic unit formatting.
    \item \texttt{\textbackslash si\{\textbackslash degreeCelsius\}} provides
          proper typesetting of the degree Celsius symbol.
    \item \texttt{\textbackslash Cref\{\}} from \ctanpackage{cleveref}
          automatically prepends ``Figure'' or ``Table'' to cross-references.
    \item The abstract environment is provided by KOMA-Script and styled
          automatically by the OmniLaTeX class.
\end{itemize}

\section{Building}

OmniLaTeX provides a Python-based build orchestrator (\texttt{build.py}) that
replaces traditional Makefiles.  To build the example article:

\begin{minted}{bash}
python3 build.py build-example article --mode dev
\end{minted}

This compiles the article through \texttt{latexmk} with three passes (sufficient
for cross-references in development mode) and places the output PDF in
\texttt{build/examples/article/main.pdf}.

To build from within Docker:

\begin{minted}{bash}
docker run --rm -v "$(pwd):/workspace" \
    ghcr.io/wyattau/omnilatex-docker:latest \
    python3 build.py build-example article --mode dev
\end{minted}

\section{Next Steps}

With the basics covered, the following chapters provide deeper dives into
specific topics:

\begin{itemize}
    \item \Cref{ch:build-system} -- Build commands, modes, caching, and
          performance optimisation.
    \item \Cref{ch:class-options} -- All class options and their effects.
    \item \Cref{ch:doctype-system} -- The full doctype profile system.
    \item \Cref{ch:doctype-system} -- Detailed guides for each document type
          (articles, theses, presentations, posters, etc.).
    \item \Cref{ch:fonts} -- Fonts, text formatting, mathematics,
          and CJK typesetting.
\end{itemize}
